\subsection{What the opened literature proves}

The sources below were opened on 18 November 2026. Page numbers refer to the
linked author version when this is stated explicitly; the bibliography records
the publication data separately. This search includes the later papers of
Hoshi--Tsujimura, Sa\"idi, Betts--Stix, and Betts--Kumpitsch--L\"udtke.
None of the scope comparisons below uses the attached unrefereed local theorem.

\begin{result}{P.1}{Known}{Nilpotent obstructions and the real case}
For the actual Galois action on the geometric fundamental group of the tripod,
the first obstruction is the integral cup product
\(\delta_2(A,B)=A\cup B\). The next obstruction is a pair of triple
Massey expressions with cyclotomic and action corrections, evaluated using
compatible defining systems belonging to one chosen class-two lift.
These are \citet[\S12.3.1, Proposition 7, author p.~14;
Theorem 19 and Remark 20(ii), author pp.~21--22]{CohWickelgren2012}.
There are non-point, noncuspidal rational abelian pairs with class-two
sections already over \(\mathbb Q\), for example \((3,4)\)
\citep[Propositions 24--25, author pp.~26--27]{CohWickelgren2012}.
For arbitrary depth, the Massey-product necessary conditions in
\citet[\S3.11 and Corollary 3.12, arXiv author pp.~16--17]{LitWickelgrenN2012}
are stated with the coefficient primes dividing \(n!\) omitted.
They are not an unrestricted integral all-depth converse at 2 and 3.
The lower-central obstruction framework also appears in the opened
\citet[pp.~1--2]{LitEllenberg2007} lecture notes. Those notes are a
2007 source, not the unpublished 2000 manuscript cited by later authors.

Over \(\mathbb R\), the full profinite section classes of a hyperbolic
curve are its real connected components
\citep[Corollary 3.13, pp.~142--143]{LitMochizuki2003}.
For a based real curve whose normalization components have real points,
the abelian classes lifting to class two already identify these components
\citep[Theorem 1.1, author p.~2]{LitWickelgrenReal2014}.
\end{result}
\begin{proof}[Derivation for this problem]
The central quotients in these obstruction statements are precisely the
geometric lower-central quotients, with the base Galois group retained.
Thus their boundary maps apply to the quotients used here. However, a class
in a boundary map has as input a chosen preceding section; vanishing of
separately chosen Massey products does not supply the compatible preceding
section required by the main question. The necessary-condition statement
with primes dividing \(n!\) excluded cannot be applied at such primes by
silently changing its coefficient ring.

For the displayed class-two example, bilinearity and the point pair
\((3,-2)\) give \(\delta_2(3,(-2)^2)=0\).
The abelian pair \((3,4)\) is neither an ordinary rational-point pair nor
one of the cusp patterns \((u,1),(1,u),(u,-u)\). It is also not an
ordinary completed pair: local full Kummer injectivity would force the
first coordinate to be the local number 3 and its complement to be both
\(-2\) and 4. Hence this counterexample does not answer the question
about \(\mathcal G_{\rm ord}\).

For the tripod over \(\mathbb R\), the three components are
\(( -\infty,0),(0,1),(1,\infty)\), with sign pairs
\((-,+),(+,+),(+,-)\). Their different abelian classes distinguish
them. In particular the real section of a full global section with rational
abelian pair is the component of that rational point. This is the real
input used in A.4, not a local theorem over \(\mathbb Q_p\).
\end{proof}

\begin{result}{P.2}{Known}{The hypotheses of the nongeometric pro-prime results}
\citet[Theorem A, author pp.~3--4; Lemma 3.2, author p.~14]{LitHoshi2010}
constructs nongeometric geometrically pro-\(p\) sections, for odd regular
\(p\), on suitable geometric \(p\)-power covers of the tripod and their
proper genus-at-least-two compactifications, after a finite extension
inside the extension unramified outside \(p\) that is pro-\(p\) over
\(\mathbb Q(\mu_p)\). The hypotheses include rational points and good
reduction away from \(p\). This theorem does not assert the claimed
construction simultaneously at every coefficient prime over \(\mathbb Q\).

\citet[Theorem 2.4 and Corollary 2.5(i), article pp.~9--11]{LitHoshiTsujimura2021}
gives a more explicit lifting statement for the tripod itself: let \(p\)
be odd regular, let \(F\) be a finite extension of \(\mathbb Q\) or
\(\mathbb Q_p\) containing \(\mu_p\), and let \(x_1,x_2\) be
\(p\)-units in \(F\cap\Omega_{p\text{-unr}}\). Then their abelian
pro-\(p\) Kummer section lifts to a full geometrically pro-\(p\)
section. The infinity-cusp qualification in P.3 is needed before using
that preprint's asserted nongeometricity criterion.

\citet[condition $(*)$ and Theorem A, \S0, pp.~573--575]{LitSaidi2023}
constructs nongeometric sections of certain geometrically
pronilpotent-by-abelian quotients of proper curves over \(p\)-adic
fields, after finite base extension. This is not a theorem about the full
pronilpotent quotient over \(\mathbb Q\).
\end{result}
\begin{proof}[Comparison of hypotheses]
For Hoshi's theorem, changing \(\mathbb Q(\mu_p)\) to \(\mathbb Q\),
changing the proper curve to the tripod, or changing one coefficient prime
to every prime is an additional assertion. P.4 below proves a particular
\(\mathbb Q\)-version directly, but still does not give ordinary pairs.
The later lifting theorem only controls ramification at the fixed finite
set consisting of \(p\). If such a pro-\(p\) abelian class were the
\(p\)-component of an ordinary completed ghost, its integral valuations
at the infinitely many support primes would be nonzero in \(\mathbb Z_p\)
and would ramify it there. This contradicts that ramification bound; see A.3.

In Sa\"idi's notation, \(\Delta\) is a pronilpotent characteristic
quotient of the geometric group, \(\widetilde H\) is its kernel,
and \(\Gamma=\Pi/\overline{[\widetilde H,\widetilde H]}\).
Condition $(*)$(iv) requires that
\(\Gamma^{(\ell)}\to\Delta^{(\ell)}\) fail to be an isomorphism
for some \(\ell\ne p\). If \(\Delta\) is the full pronilpotent
quotient, its pro-\(\ell\) quotient is already the full
\(\Pi^{(\ell)}\). Every quotient between \(\Pi\) and \(\Delta\),
including \(\Gamma\), has that same pro-\(\ell\) quotient.
Condition $(*)$(iv) then fails. Thus even removing the base-field and
properness differences would not make this theorem an all-prime limit theorem.
\end{proof}

\begin{result}{P.3}{Refuted}{The infinity cusp must be included in the 2021 criterion}
The implication from condition (3) to ``(2) or (4)'' in
\citet[Proposition 1.4(ii), article p.~6]{LitHoshiTsujimura2021}
omits the cusp at infinity and is false as printed.
Consequently Theorem A(iii) and Corollary 2.5(iii) cannot be used without
that qualification. The definition invoked there is
\citet[Definition 1.1(iii), author p.~7]{LitHoshi2010}: geometricity means
containment in a decomposition group of a rational point of the smooth
compactification. The 2021 Theorem A(iii), article p.~3, repeats this
definition explicitly; it includes the point at infinity.

Here is the corrected tripod statement. For a number field or a finite
extension \(F/\mathbb Q_p\), write \(D_p\subset F^\times\) for
the kernel of the pro-\(p\) Kummer map; it consists of the roots of unity
of order prime to \(p\). A pair of actual elements
\((x_1,x_2)\in(F^\times)^2\) is the abelian image of at least one
geometric geometrically pro-\(p\) section, where geometric includes
cusps, if and only if at least one of the following holds:
\[
\begin{array}{ll}
\text{ordinary:}&\zeta_1x_1+\zeta_2x_2=1
                  \quad\text{for some }\zeta_1,\zeta_2\in D_p;\\
\text{cusp }0:&x_2\in D_p;\\
\text{cusp }1:&x_1\in D_p;\\
\text{cusp }\infty:&-x_2/x_1\in D_p.
\end{array}
\]
Thus the lift in P.2 is certainly nongeometric if all four conditions fail.
\end{result}
\begin{proof}
The ordinary condition is equality of the two Kummer coordinates with
\((t,1-t)\). At the three rational cusps, the images of the decomposition
sections in the abelian section set have, respectively, the forms
\((u,1),(1,u),(u,-u)\), with \(u\) in the pro-\(p\) completion.
For the last formula put \(z=1/t\). Then
\(t=z^{-1}\) and \(1-t=-z^{-1}(1-z)\); a tangent with
\(z\)-coordinate \(v\ne0\) gives \((v^{-1},-v^{-1})\).
General decomposition sections differ by inertia Kummer classes and
give the same three patterns with completed \(u\).
These calculations prove necessity. For actual \(x_i\), every listed
condition is sufficient: choose the indicated ordinary point, or choose
the tangent having the required actual coordinate \(u\). The kernel
description converts equality of completed coordinates to the displayed
conditions. This proves the corrected equivalence.

For a counterexample satisfying the 2021 preprint's hypotheses take
\(p=3\), \(F=\mathbb Q(\mu_3)\), and \(x_1=x_2=3\).
The prime-to-3 roots of unity in \(F\) are \(\{1,-1\}\), and
\(3\zeta_1+3\zeta_2\in\{-6,0,6\}\), so its ordinary criterion
fails. Neither coordinate is a root of unity. Nevertheless the tangent
\(v=1/3\) at infinity supplies a cuspidal section whose pair is
\((3,-3)=(3,3)\) in pro-3 Kummer classes: \(-1\) is a
\(3^n\)-th power for every \(n\). This contradicts the asserted
nongeometricity. The finite root-of-unity calculation is checked in
\texttt{code/check\_literature\_scope.py}; the tangent calculation proves
the full inverse-limit assertion, rather than inferring it from that
finite check. The lifting assertion itself is unaffected.
\end{proof}

\begin{result}{P.4}{Proved}{Full pro-prime counterexamples over \(\mathbb Q\), including 2}
Let \(p\) be 2 or an odd regular prime. There is a section of the full
geometrically pro-\(p\) fundamental extension of the tripod over
\(\mathbb Q\) with abelian pair
\((\kappa_p(p),\kappa_p(p^2))\). Every such section is nongeometric,
even allowing cusps. Its abelian pair is not ordinary at \(p\), so it is
not a counterexample to either limit assertion for ordinary completed pairs.
For odd regular \(p\), the same construction lifts every pair
\((a\kappa_p(p),b\kappa_p(p))\), \(a,b\in\mathbb Z_p\).
This is a derivation from the cited lifting and Galois-group theorems, not
a claim that the quoted papers state this precise \(\mathbb Q\)-form.
\end{result}
\begin{proof}
Choose the rational tangential base point \(0+1\), write \(P\) for the
free pro-\(p\) geometric group, and identify its arithmetic extension
with \(P\rtimes G_{\mathbb Q}\). Define \(\Omega_{p\text{-unr}}\)
to be the maximal pro-\(p\) extension of \(\mathbb Q(\mu_p)\)
unramified outside \(p\), inside the chosen algebraic closure. It is
Galois over \(\mathbb Q\). For \(p=2\) this means the maximal
pro-2 extension of \(\mathbb Q\) unramified outside 2, with ramification
at the real place allowed. This is not the unrestricted maximal
extension of \(\mathbb Q\) unramified outside \(p\).
The pointed action on \(P\) is
unramified at every finite prime different from \(p\): the tripod and
its boundary have smooth good reduction there, and the tangent has unit
coordinate. This is the pointed specialization construction, equivalently
the argument of \citet[Lemma 2.2 and its proof, article pp.~8--9]{LitHoshiTsujimura2021}
applied after adjoining \(\mu_p\). On the Frattini quotient
\(P/\Phi(P)\cong\mathbb F_p(1)^2\), the action is scalar
cyclotomic. The kernel of
\(\operatorname{Aut}(P)\to\operatorname{Aut}(P/\Phi(P))\)
is pro-\(p\). Therefore the pointed action factors through
\[
 \Gamma_p=\operatorname{Gal}(\Omega_{p\text{-unr}}/\mathbb Q).
\]
The Kummer cocycle \(\kappa_p(p)\) also factors through this quotient:
its defining radical and root-of-unity extensions are unramified away
from \(p\) and pro-\(p\) after adjoining \(\mu_p\).

Suppose first that \(p\) is odd regular. The subgroup
\(H=\operatorname{Gal}(\Omega_{p\text{-unr}}/\mathbb Q(\mu_p))\)
is free pro-\(p\) by \citet[Lemma 3.2, author p.~14]{LitHoshi2010},
and \([\Gamma_p:H]=p-1\). Hochschild--Serre, or restriction and
corestriction, gives \(\operatorname{cd}_p\Gamma_p\le1\).
Put \(M=P^{\rm ab}\). The graph of any displayed abelian cocycle is
a continuous homomorphism \(\Gamma_p\to M\rtimes\Gamma_p\).
Apply the pro-\(p\)-kernel lifting theorem to
\[
 P\rtimes\Gamma_p\longrightarrow M\rtimes\Gamma_p.
\]
Its kernel is \(\overline{[P,P]}\), a pro-\(p\) group. The required
lifting criterion is explicitly stated in
\citet[Proposition 4.3, p.~149]{LitPachecoStevenson2000}, citing
Serre's cohomological-dimension theorem. A lifted homomorphism still
projects identically to \(\Gamma_p\). Its crossed homomorphism,
inflated to \(G_{\mathbb Q}\), is the required section. No descent
of an arbitrarily chosen section over \(\mathbb Q(\mu_p)\) is asserted;
the lifting problem is solved over \(\Gamma_p\) itself.

For \(p=2\), the field \(\mathbb Q\) is 2-rational by criterion
(iv) of \citet[Theorem and Definition 0, author p.~2]{LitGrasJaulent2009}.
Indeed a rational number with even valuations at every odd prime is,
modulo rational squares, one of \(1,-1,2,-2\), and only 1 is a
square in \(\mathbb Q_2\). The global and local 2-primary
root-of-unity ranks both equal one. Their Theorem 2, author p.~3, gives
\(\Gamma_2\cong\mathbb Z_2*C_2\), with the \(C_2\) factor
the real decomposition group. Choose compatible positive real roots of 2.
The abelian cocycle \(\alpha=(\kappa_2(2),2\kappa_2(2))\) is then
zero on complex conjugation. On the \(C_2\) factor lift its graph by
the reference tangential section. On the free \(\mathbb Z_2\) factor
choose any lift of the value of its generator through
\(P\rtimes\Gamma_2\to M\rtimes\Gamma_2\).
Since the source group is pro-2, powers of this element define a continuous
homomorphism from \(\mathbb Z_2\). The universal property of the
pro-2 free product combines the two lifts. Their abelianization equals
the prescribed graph on both factors and hence on \(\Gamma_2\).
Inflating again yields a full geometrically pro-2 section over \(\mathbb Q\).

Finally the two valuations of the pair at \(p\) are \((1,2)\).
An ordinary local point with \(v_p(t)>0\) has \(v_p(1-t)=0\).
The cusp valuation patterns are \((m,0),(0,m),(m,m)\), even when
\(m\) is allowed in \(\mathbb Z_p\). The pair \((1,2)\) has none
of these forms. Localization therefore excludes every geometric section
and proves the asserted failure of local ordinarity.
\end{proof}

\begin{result}{P.5}{Known}{Unipotent Selmer results and the missing hypotheses}
Kim's theorem for the tripod proves finiteness of
\(U(\mathbb Z[S^{-1}])\) for a fixed finite set \(S\), using
rational unipotent fundamental groups and Selmer varieties
\citep[\S3, concluding proof of Siegel's theorem, author pp.~17--18]{LitKim2005}.
The conjecture of \citet[\S1.3 and Conjecture 3.1,
author pp.~2--3 and 12; finite-$S$ variant in \S8]{LitBDCKW2018}
asserts eventual equality of the associated local Chabauty--Kim locus
with the integral points of the chosen model. These Selmer classes have
fixed ramification and the stated local geometric and crystalline
conditions.

The strengthened comparison in
\citet[Theorem A; Definition 2.1; \S2.2.1 and Lemma 2.19]{LitBKL2026}
still concerns a fixed finite \(S\) and adelic, locally geometric
sections. Its Theorem B and Corollary 5.17 establish the refined
Chabauty--Kim conjecture for the tripod over \(\mathbb Z[1/2]\),
for every odd coefficient prime. Its Remark 1.8 explicitly records that
the corresponding full Selmer section assertion over \(\mathbb Q\),
for the tripod and any finite \(S\), was already a consequence of
\citet[Corollary 6, author p.~6]{Stix2015}.
For the tripod there is a useful stronger local comparison:
\citet[Lemma 4.2 and Corollary 4.3, author pp.~31--32]{LitBKL2026}
shows that away from the coefficient prime \(p\), rational unipotent
local cohomology is already determined by abelianization.
\end{result}
\begin{proof}[Precise comparison with the main question]
The Malcev completion homomorphism
\(P_p\to U_{p,n}(\mathbb Q_p)\) is Galois equivariant, and it sends
path torsors to their unipotent path torsors. Thus a full pro-\(p\)
section has a well-defined rational unipotent image at every depth.
The cited Lemma 2.19 verifies compatibility with the local Kummer maps.
If the section satisfies the relevant full local point conditions and
fixed finite-ramification conditions, this image is in the relevant
Selmer set; for an integral point at \(p\), the usual crystalline
condition is also the one used there.

At a place \(q\ne p\), the last comparison does identify the unipotent
local class with that of the point supplied by ordinarity: for every
\(j\ge2\), both \(H^1(G_q,\mathbb Q_p(j))\) and
\(H^2(G_q,\mathbb Q_p(j))\) vanish, and the central tower gives a
bijection on first cohomology. At \(q=p\), P.9 below shows that all
these torsors are semistable, and that a torsor with integral ordinary
abelianization is crystalline. Thus higher Hodge admissibility is not
an additional missing condition here. Equality of abelian classes still
does not specify the higher local path coordinates. The main hypothesis
therefore does not give the full local Kummer equality at the coefficient
place that is needed to place its specified point in a Chabauty--Kim locus.
Moreover an infinite-support ordinary pair is not integral outside any
fixed finite \(S\), and by A.3 its \(p\)-component ramifies at
infinitely many primes. It therefore fails an independent hypothesis of
every fixed-\(S\) Selmer construction. Taking a union over finite sets
does not include an element of infinite support. Finally rational
unipotent coefficients discard integral and torsion information, so the
forward map supplies no reverse lift of a unipotent Selmer class to an
integral pro-\(p\), much less all-prime pronilpotent, section.
These are the precise gaps in an attempted direct application of Kim's
theorem or conjecture to the question in A.2.
\end{proof}

\begin{result}{P.6}{Known}{Continuous cohomology and the scope of a derived-limit obstruction}
For \(T=\varprojlim_m T_m\), with finite discrete coefficient modules
and surjective transition maps, continuous cohomology satisfies
\[
0\longrightarrow\varprojlim{}^1_m H^{i-1}(G,T_m)
\longrightarrow H^i_{\rm cts}(G,T)
\longrightarrow\varprojlim_m H^i(G,T_m)\longrightarrow0.
\]
This is \citet[equation (2.1) and Theorem (2.2), p.~215;
Proposition (1.6), p.~211]{CohJannsen1988}.
Finite cohomology groups at a local field, or with a fixed finite set of
global ramification primes, make the relevant tower Mittag--Leffler and
kill its \(\varprojlim^1\). Finiteness of the modules \(T_m\)
alone does not imply finiteness of \(H^1(G_{\mathbb Q},T_m)\).
In particular the unrestricted weight-two and weight-twelve computations
in Part A require their separate arithmetic inputs and do not follow
from Jannsen's formalism alone.
\end{result}
\begin{proof}[Derivation and limit of the conclusion]
Take \(T_m=\mathbb Z/p^m(j)\) and \(i=2\).
The left-hand term consists of continuous classes whose image at every
finite coefficient level is zero. At each completion, the inverse system
of finite first-cohomology groups is Mittag--Leffler, so the local
derived-limit term vanishes. Naturality then makes every such global
derived-limit class locally trivial. Globally, already the weight-one
finite-coefficient Kummer groups allow arbitrary finite sets of ramified
primes, so they are infinite; one cannot reuse the fixed-ramification
Mittag--Leffler argument for the unrestricted group.

This statement identifies a possible target for an obstruction. It does
not show that a specified ordinary pair, or one of its specified lifts,
has a nonzero obstruction in that target. The arithmetic calculation of
the weight-two restricted-product quotient, the weight-twelve/691 class,
and their relation to actual obstruction cochains is consequently kept
separate in Part A. In particular existence of a nonzero localization
kernel does not prove Lead L4.
\end{proof}

\begin{result}{P.7}{Known}{Stronger full-section theorems do not remove the quotient gap}
For a hyperbolic curve over \(\mathbb Q\) with a nonconstant unit,
\citet[Corollary 6, author p.~6]{Stix2015} identifies its full adelic
section classes with rational points; Corollary 5, author p.~5, also
handles noncuspidal Selmer sections having finite geometric support.
For a smooth proper genus-at-least-two curve over a number field containing
no CM subfield, \citet[Theorem A and Definition 1.1,
author p.~1]{LitBettsStix2025} proves finiteness of the image at every
finite place of the full Selmer section set.
Neither statement is a theorem about arbitrary pronilpotent sections
whose abelian localizations are ordinary.
\end{result}
\begin{proof}[Reduction relevant here]
The tripod has the nonconstant unit \(t\), so the first theorem applies
as soon as full local point sections and adelicity have been established.
A rational abelian pair makes all finite local points equal to the
rational point, by local Kummer injectivity; hence it supplies those
remaining support and adelicity conditions. This is exactly the
derivation in A.4--A.5, which strengthens the implication claimed in the
prompt. For a section only of the pronilpotent quotient, the reverse
lift to the full fundamental group is an additional embedding problem.
Betts--Stix also starts with full local point sections, and uses
representations furnished by abelian-by-finite families; no theorem there
states that these representations factor through the pronilpotent
quotient of the tripod. Thus its finiteness result does not supply the
missing hypothesis either.
\end{proof}

\begin{result}{P.8}{Proved}{What this literature comparison says about R1--R7}
The cup-product and Massey descriptions, the formal lower-central
obstruction tower, real component detection, the fixed-finite-\(S\)
unipotent framework, and the derived-limit exact sequence are established
inputs and are not new results of the earlier runs. The specific
factorial-prime ordinary-pair constructions in R1 and R3, their
finite-prefix exclusions in R5, the pointed Heisenberg-cover arithmetic
in R6, and the uniform exclusion assertion in R7 are not statements
proved in the sources opened here. Their validity is determined by the
independent audits in Parts A--E, not by attribution to those sources.
No priority claim for the earlier constructions is inferred from this
bounded literature search.
\end{result}
\begin{proof}
P.1--P.7 and P.9--P.11 specify the coefficient groups, local conditions, ramification
conditions, and field hypotheses of each source used. None is an
existence theorem for an infinite-support ordinary completed pair with
a section of the full pronilpotent quotient, nor a finite-support theorem
under precisely that hypothesis. This is a comparison of the opened
statements, not a claim that every possible source has been exhausted.
The independent proofs and refutations elsewhere in this paper are needed
before any of R1--R7 can be used as a mathematical premise.
\end{proof}

\begin{result}{P.9}{Proved}{The precise local Hodge bridge for the tripod}
Let \(U_{p,n}\) be the depth-\(n\) rational \(\mathbb Q_p\)-unipotent
geometric fundamental group of the tripod, based at \(0+1\), with its
continuous \(G_{\mathbb Q_p}\)-action. Every class of
\(H^1(G_{\mathbb Q_p},U_{p,n}(\mathbb Q_p))\) is semistable and
hence de Rham. Moreover it is crystalline if and only if its abelian
image is crystalline. Consequently, if a pro-\(p\) section has ordinary
abelian image represented by \(x_p\), its rational unipotent image is
semistable at every depth, and is crystalline when
\(x_p,1-x_p\in\mathbb Z_p^\times\).
These conclusions do not assert equality with the entire path torsor
of \(x_p\).
\end{result}
\begin{proof}
The weight filtration on \(\mathcal O(U_{p,n})\) has finite-dimensional
pieces and graded pieces that are sums of \(\mathbb Q_p(-j)\),
\(j\ge0\), in weight \(2j\). Translation by an element of
\(U_{p,n}\) preserves this filtration and acts identically on its
associated graded. Identify an algebraic \(U_{p,n}\)-torsor with
\(U_{p,n}\) using a \(\mathbb Q_p\)-point, which exists because a
unipotent group over a characteristic-zero field has no nontrivial
algebraic torsors. Twisting the Galois action by its continuous cocycle
therefore gives the same graded Tate pieces on its coordinate algebra.
Each finite weight piece is a mixed Tate representation in the precise
sense of \citet[Definition 1.3.7]{LitChatzistamatiouUnver2013}.
Their Proposition 1.3.8, arXiv author pp.~12--13, proves that every such
representation is semistable. Applying the exact tensor functor
\(D_{\rm st}\) to these pieces produces a torsor under
\(D_{\rm st}(U_{p,n})\) over \(\mathbb Q_p\); this unipotent
torsor again has a rational point. The comparison isomorphism gives
a Galois-invariant \(B_{\rm st}\)-point of the original torsor.
This is the semistable condition. This argument applies also at 2 and
does not require integrality or even ordinary abelianization.

For crystallinity use the central tower. Its successive kernels are
\(V_j=\mathbb Q_p(j)^{b_j}\), \(j\ge2\).
\citet[Appendix A, author p.~47, citing Bloch--Kato, Example 3.9]{LitBKL2026}
records \(H_f^1(G_{\mathbb Q_p},V_j)=H^1(G_{\mathbb Q_p},V_j)\).
The natural nonabelian Bloch--Kato isomorphism
\[
 H_f^1(G_{\mathbb Q_p},U_{p,j})
       \simeq U_{p,j}^{\rm dR}(\mathbb Q_p)
\]
in \citet[\S3.3, equation (3.8), author p.~27]{LitBKL2026}
commutes with the depth projections. These projections on de Rham groups
are surjective on \(\mathbb Q_p\)-points, so every crystalline class
at depth \(j-1\) has a crystalline lift to depth \(j\).
The fibre of the unrestricted nonabelian cohomology map over that class
is a transitive \(H^1(G_{\mathbb Q_p},V_j)\)-set.
Twisting a crystalline torsor by a crystalline central torsor preserves
crystallinity: over \(B_{\rm cris}\), multiply their two trivializing
points, using centrality. Since every such central twisting class is
crystalline, every lift is crystalline. Induction proves the converse
direction; the forward direction follows by taking the abelian quotient.
At depth one the crystalline Kummer subspace is the image of units,
which proves the stated consequence for \(x_p\).

There remain higher coordinates even in the crystalline subset. Already
at depth two the fibre over a fixed crystalline abelian class is a
torsor under the one-dimensional group
\(H^1(G_{\mathbb Q_p},\mathbb Q_p(2))\); it is free because the
abelian quotient has no Galois invariants. Thus specifying the
abelianization and Hodge admissibility does not specify a path class.
In particular this result removes a potential Hodge gap in P.5 but
neither its local-point-equality gap at \(p\) nor its independent
fixed-finite-\(S\) gap.
\end{proof}

\begin{result}{P.10}{Known}{Brauer--Manin and descent comparisons}
\citet[Theorem 1.2, corrected author pp.~2--3; Remark 2.4, p.~7]{LitHarariSzamuely2009}
characterizes the splitting of the Albanese quotient of the geometric
abelian fundamental extension by membership of the Albanese torsor in
the maximal divisible subgroup of \(H^1(k,\operatorname{Alb}_X)\).
For curves the geometric abelianization itself is the relevant Tate
module. The quasi-projective version uses the generalized Albanese.
For a torsor under a semiabelian variety, finite Tate--Shafarevich group
of its abelian quotient makes the indicated Brauer--Manin obstruction
complete for existence of a rational point
\citep[Theorem 1.1, author p.~2; Corollary 4.8 and Lemma 4.9,
author p.~16]{LitHarariSzamuely2008}.

For a geometrically connected variety over a number field and a continuous
quotient of its geometric fundamental group whose kernel is normal in the
arithmetic fundamental group,
\citet[Theorem 2.1 and Remark (3), author pp.~11--12]{LitHarariStix2012}
equates interpolation of a tuple of local point sections in that quotient
with survival of the corresponding finite descent obstructions.
This includes the full pronilpotent quotient and does not require an
adelic tuple when formulated as that interpolation statement.
Their Corollary 3.1, author p.~13, gives Brauer--Manin orthogonality
on a proper curve from a section of its full geometric abelianization
whose local sections come from points.
\end{result}
\begin{proof}[Why these theorems do not yet give the required rigid target]
For the tripod, the generalized Albanese is the split two-dimensional
torus with immersion \(t\mapsto(t,1-t)\), and its Albanese torsor is
already trivial. The splitting-existence criterion therefore adds no
restriction to the chosen abelian section class. Likewise a rational
point on this torus, or density of its rational points in a product of
local tori, does not make the specified local tuple a single rational
point of the immersed curve. The semiabelian torsor theorem is an
existence statement for a different variety, not this stronger assertion.

To apply Harari--Stix with the pronilpotent quotient, one must show
that its local sections equal the local point sections in that quotient.
The hypothesis defining \(\mathcal G_{\rm ord}\) gives this only
in the abelian quotient. P.5 and P.9 give useful rational unipotent
comparisons, but not the required integral pronilpotent equality.
Finally, passing to a finite nilpotent cover does not automatically
provide a section of its full geometric abelianization. The image of
the covering subgroup in the base's pronilpotent quotient need not
retain that subgroup's full Tate modules at all primes. In particular
a finite-index pro-\(p\) subgroup retains its pro-\(p\) completion,
but its additional abelian quotients at primes different from \(p\)
need not be present in the base's pronilpotent quotient. A cover-based
application must establish that comparison, as well as local point
conditions, before the cited Brauer--Manin conclusion transfers.
No such bridge for every ordinary completed pair is supplied by these
theorems.
\end{proof}

\begin{result}{P.11}{Known}{Completed reciprocity complexes and their point domain}
\citet[Definitions 3.1--3.3, p.~728; Definition 3.10, p.~733]{LitPridham2020}
defines adelic cohomology for profinite coefficients by a derived inverse
limit of finite-coefficient restricted products, and compactly supported
cohomology by the cocone of global restriction. These constructions allow
an infinite set of finite places. The nilpotent reciprocity tower in
Example 3.16, pp.~734--735, nevertheless starts from actual adelic points
and requires a global quotient section matching their local path classes.
Example 3.18, pp.~735--736, gives the rational unipotent variant.
\end{result}
\begin{proof}[Consequences for completed ordinary pairs]
The two restrictions must be distinguished. Every continuous
finite-coefficient global cohomology class is unramified outside finitely
many places: its cocycle and the coefficient action factor through a
finite Galois extension, which is ramified at finitely many primes.
The finite set may grow as the coefficient quotient grows. Thus the
abelian localizations of a completed ordinary pair lie in the
cohomological restricted product even when the underlying ordinary
points have infinite support. Infinite support alone does not invalidate
this cohomological construction.

For a profinite Tate module \(M\), the defining triangle gives
\[
 H^1(G_{\mathbb Q},M)\xrightarrow{\operatorname{res}}
 H^1(G_{\mathbb A_f},M)\xrightarrow{\partial}
 H_c^2(G_{\mathbb Q},M).
\]
Consequently \(\partial\operatorname{res}(A)=
\partial\operatorname{res}(B)=0\) for a completed global Kummer pair.
This is an unconditional formal consequence of the cocone, independent
of a duality identification. It proves that this first abelian boundary
cannot by itself exclude an ordinary ghost.

In contrast, a point of the standard tripod over the finite adeles has
both coordinates units at almost every prime. An ordinary tuple with
infinite valuation support is therefore outside that point domain.
Even after enlarging the domain formally, a higher point reciprocity
condition would require compatibility between the chosen higher local
point paths and the global lift, whereas ordinary abelianization supplies
only depth one. The existence of a global pronilpotent lift does not
identify its localization with those point paths. The cohomological
completion and an infinite set of allowed primes remove neither issue.

No Poitou--Tate identification omitting real-place corrections is used
here at 2. A duality interpretation over \(\mathbb Q\) must include the
appropriate real Tate terms; the formal cone and the displayed vanishing
remain valid without that additional identification.
\end{proof}
